\begin{figure}[H]
	\centering 
	\includegraphics[width=1.0\textwidth]{SpecificaTecnica/Recording/Diagramma}
	\caption{Diagramma delle classi per Recording}
\end{figure}

\subsubsubsection{RecordingController}
\begin{figure}[H]
	\centering 
	\includegraphics[width=0.6\textwidth]{SpecificaTecnica/Recording/RecordingController}
	\caption{Recording: RecordingController}
\end{figure}

\paragraph{Descrizione}
Modulo FastAPI che definisce il router per le API$^G$ REST$^G$ della trascrizione audio.\\
Espone un unico endpoint per la ricezione di file audio e la restituzione della trascrizione testuale.\\
Delega la logica applicativa a \textit{RecordingService}, ottenuto tramite dependency injection insieme al client OpenAI.

\paragraph{Attributi}
\begin{itemize}
	\item \textbf{router: APIRouter} - router FastAPI con prefisso \textit{/recording} che raggruppa gli endpoint del modulo.
	\item \textbf{ServiceDep} - tipo annotato che inietta un'istanza di \textit{RecordingService} tramite la funzione \textit{get\_service}.
\end{itemize}

\paragraph{Metodi e funzioni}
\begin{itemize}
	\item \textbf{get\_service(): RecordingService} - funzione di dependency injection che istanzia \textit{AsyncOpenAI}, \textit{RecordingRepository}, \textit{RecordingAdapter} e \textit{RecordingService}.
	\item \textbf{trascrivi\_audio(file: UploadFile, service: RecordingService): Dict[string, string]} - endpoint \textit{POST /recording/trascrivi} che riceve un file audio e restituisce la trascrizione testuale nel campo \textit{testo}.\\
	Lancia \textit{HTTPException} con codice 400 se il file è vuoto, con codice 500 se la trascrizione fallisce.
\end{itemize}

\subsubsubsection{RecordingService}
\begin{figure}[H]
	\centering 
	\includegraphics[width=0.6\textwidth]{SpecificaTecnica/Recording/RecordingService}
	\caption{Recording: RecordingService}
\end{figure}

\paragraph{Descrizione}
Classe$^G$ che implementa la logica applicativa per la gestione della trascrizione audio.\\
Prima di invocare la trascrizione, valida il file ricevuto verificando che non sia vuoto, che non superi i 10 MB e che la durata non ecceda i 120 secondi.\\
La validazione$^G$ della durata viene saltata per i formati non supportati da \textit{mutagen}, come webm/opus.\\
Delega la trascrizione effettiva all'istanza di \textit{RecordingPort} ricevuta alla costruzione.

\paragraph{Attributi}
\begin{itemize}
	\item \textbf{adapter: RecordingPort} - istanza della porta di trascrizione iniettata alla costruzione.
	\item \textbf{MAX\_DURATION\_SEC: int} - costante che definisce la durata massima consentita per il file audio, impostata a 120 secondi.
	\item \textbf{MAX\_SIZE\_BYTES: int} - costante che definisce la dimensione massima consentita per il file audio, impostata a 10 MB.
\end{itemize}

\paragraph{Metodi e funzioni}
\begin{itemize}
	\item \textbf{\_\_init\_\_(adapter: RecordingPort)} - inizializza il service con l'istanza della porta di trascrizione ricevuta.
	\item \textbf{trascrivi\_audio(audio\_bytes: bytes, filename: string): string} - valida il file audio e, se i controlli sono superati, invoca \textit{adapter.trascrivi} per ottenere il testo trascritto.\\
	Lancia \textit{ValueError} se il file è vuoto, supera i 10 MB o supera i 120 secondi di durata.
\end{itemize}

\subsubsubsection{RecordingPort}
\begin{figure}[H]
	\centering 
	
\includegraphics[width=0.6\textwidth]{SpecificaTecnica/Recording/RecordingPort}
	\caption{Recording: RecordingPort}
\end{figure}

\paragraph{Descrizione}
Classe astratta$^G$ che definisce la porta per la trascrizione audio tramite servizio esterno.\\
Segue il pattern architetturale delle porte e degli adattatori, disaccoppiando la logica applicativa di \textit{RecordingService} dalla specifica implementazione del servizio di trascrizione.

\paragraph{Attributi}
Essendo un'interfaccia, non definisce attributi propri.

\paragraph{Metodi e funzioni}
\begin{itemize}
	\item \textbf{trascrivi(audio\_bytes: bytes, filename: string): string} - trascrive un file audio in testo a partire dal contenuto binario e dal nome del file.\\
	Lancia un'eccezione se la trascrizione fallisce.\\
\end{itemize}

\subsubsubsection{RecordingAdapter}
\begin{figure}[H]
	\centering 
	\includegraphics[width=0.5\textwidth]{SpecificaTecnica/Recording/RecordingAdapter}
	\caption{Recording: RecordingAdapter}
\end{figure}

\paragraph{Descrizione}
Classe$^G$ che implementa \textit{RecordingPort} seguendo il pattern degli adattatori.\\
Riceve un'istanza di \textit{RecordingRepository} alla costruzione e delega interamente la trascrizione audio.

\paragraph{Attributi}
Oltre all'implementazione concreta della funzione di \textit{RecordingPort}, è disponibile la seguente funzione:
\begin{itemize}
	\item \textbf{\_\_init\_\_(repo: RecordingRepository)} - inizializza l'adapter con l'istanza del repository$^G$ ricevuta.
\end{itemize}

\paragraph{Metodi e funzioni}
\begin{itemize}
	\item \textbf{trascrivi(audio\_bytes: bytes, filename: string): string} - delega a \textit{RecordingRepository} la trascrizione del file audio.
\end{itemize}

\subsubsubsection{RecordingRepository}
\begin{figure}[H]
	\centering 
	\includegraphics[width=0.6\textwidth]{SpecificaTecnica/Recording/RecordingRepository}
	\caption{Recording: RecordingRepository}
\end{figure}

\paragraph{Descrizione}
Classe$^G$ che gestisce l'accesso diretto al servizio di trascrizione audio Whisper$^G$ tramite il client \textit{AsyncOpenAI}.\\
Riceve il client alla costruzione e lo utilizza per inviare il file audio e recuperare il testo trascritto.

\paragraph{Attributi}
\begin{itemize}
	\item \textbf{client: AsyncOpenAI} - istanza del client OpenAI$^G$ asincrono utilizzata per le chiamate al modello Whisper.
\end{itemize}

\paragraph{Metodi e funzioni}
\begin{itemize}
	\item \textbf{\_\_init\_\_(client: AsyncOpenAI)} - inizializza il repository$^G$ con il client OpenAI$^G$ asincrono ricevuto.
	\item \textbf{trascrivi(audio\_bytes: bytes, filename: string): string} - invia il contenuto binario del file audio al modello Whisper$^G$ tramite il client OpenAI$^G$ e restituisce il testo trascritto.\\
	Lancia un'eccezione se la chiamata al servizio fallisce.
\end{itemize}